\subsection{Steuding's orbit continued fraction: an exact endpoint recoding}
\label{subsec:Steuding-522}

Steuding's starred Exercise~5.22 uses the unshortened positive-integer
recursion \eqref{eq:unshortened-U}, asks for an equivalence-class
reformulation of the Collatz conjecture, and asks for a Markoff constant
\cite[physical PDF p.~100; printed p.~87, Ex.~5.22]{Steuding2005}.
It is an exercise prompt, not a theorem accompanied by a proof.  The
propositions and calculation below supply an edition solution without a
novelty claim.

The target printed in the exercise is the purely periodic irrational
\[
 \beta=[\overline{4,2,1}],
\]
with a visible overbar.  Ordinary PDF text extraction loses that bar.  The
unbarred finite fraction is $[4,2,1]=13/3$ and cannot be substituted:
Steuding's rational and irrational equivalence classes are disjoint
\cite[printed p.~79]{Steuding2005}.

\begin{proposition}[whole orbit to continued-fraction morphism]
\label{prop:Steuding-orbit-CF-morphism}
For $m\in\Npos$, define
\begin{equation}
\label{eq:Steuding-Theta}
 x_j(m)=U^j(m),
 \qquad
 \Theta(m)=[x_0(m),x_1(m),x_2(m),\ldots].
\end{equation}
Every $\Theta(m)$ is an irrational real number.  If
\[
 \mathcal G(y)=\frac{1}{y-\lfloor y\rfloor}
\]
is the Gauss tail map on nonintegral reals, then
\begin{equation}
\label{eq:Steuding-Gauss-conjugacy}
 \lfloor\Theta(m)\rfloor=m,
 \qquad
 \mathcal G(\Theta(m))=\Theta(U(m)).
\end{equation}
Consequently $\Theta$ is a bijection from $\Npos$ onto its image, its inverse
there is the floor map, and it conjugates $U$ to $\mathcal G$ restricted to
that image.
\end{proposition}

\begin{proof}
Every $x_j(m)$ is a positive integer.  The finite convergents of
\eqref{eq:Steuding-Theta} have denominator continuants tending to infinity,
and consecutive convergents differ by the reciprocal of the product of
their denominators.  Their alternating nested intervals therefore have a
unique common limit.  That limit cannot be rational: if it were the reduced
fraction $p/q$, choose $n$ so large that $q_n>q$ and $q_{n+1}>q$.  Since each
convergent $p_n/q_n$ is reduced, it is then distinct from $p/q$.  The source
bound
\[
 \left|\frac{p}{q}-\frac{p_n}{q_n}\right|<
 \frac{1}{q_nq_{n+1}}
\]
would make the left-hand side strictly smaller than $1/(q q_n)$.  But two
distinct reduced rationals with
denominators $q$ and $q_n$ are separated by at least $1/(q q_n)$.  This is
the constructive
content of Steuding's infinite-fraction discussion
\cite[Sec.~3.4, printed pp.~42--43]{Steuding2005}.

The tail $[x_1(m),x_2(m),\ldots]$ is greater than one, so
\[
 \Theta(m)=m+\frac{1}{[x_1(m),x_2(m),\ldots]}
\]
has floor $m$.  Applying $\mathcal G$ deletes the first digit, while the
recursion gives $x_{j+1}(m)=x_j(U(m))$.  This proves both identities in
\eqref{eq:Steuding-Gauss-conjugacy}.  The floor identity proves injectivity
and supplies the displayed inverse; surjectivity is only onto the stated
image.
\end{proof}

Steuding defines $\alpha\sim\gamma$ by
\begin{equation}
\label{eq:Steuding-unimodular-equivalence}
 \alpha=\frac{a\gamma+b}{c\gamma+d},
 \qquad
 a,b,c,d\in\Z,
 \qquad
 ad-bc=\pm1.
\end{equation}
Serret's Theorem~5.7 says that two irrational simple continued fractions are
equivalent exactly when finite prefixes can be removed from both to leave the
same infinite tail \cite[printed pp.~79--80]{Steuding2005}.

\begin{sourceissue}[the matrix step in the printed Serret proof]
\label{sourceissue:Steuding-Serret-proof}
Printed p.~78 calls all integer matrices of determinant $\pm1$ the special
linear group $\mathrm{SL}_2(\Z)$; the group is $\mathrm{GL}_2(\Z)$, since
$\mathrm{SL}_2(\Z)$ has determinant $+1$.  Lemma~5.6 also suppresses the
choice of simultaneous matrix sign and leaves its constructive sufficiency
direction to an induction.

The converse proof on printed p.~80 then (i) invokes $c>d>0$ for the original
matrix before that reduction has been proved, (ii) duplicates the numerator
entry $R$ by printing
$S=ap_{m-1}+bq_{m-1}$ instead of
$S=cp_{m-1}+dq_{m-1}$, and (iii) writes $\alpha$ where the convergents and
error terms belong to $\gamma$ (called $\beta$ in the source).  The factors
$c\alpha+d$ on the same page must likewise be $c\gamma+d$.

Here is the exact repair.  Start with an arbitrary matrix $M$ in
\eqref{eq:Steuding-unimodular-equivalence} and replace it by $-M$, if
necessary, so that $D=c\gamma+d>0$.  For convergents $p_m/q_m$ to $\gamma$,
put $\varepsilon_m=p_m-\gamma q_m$.  The lower row of the composite matrix
is
\begin{align*}
 Q_m&=cp_m+dq_m=Dq_m+c\varepsilon_m,\\
 S_m&=cp_{m-1}+dq_{m-1}=Dq_{m-1}+c\varepsilon_{m-1}.
\end{align*}
Since $\varepsilon_m\to0$ and
\[
 q_m-q_{m-1}=(a_m-1)q_{m-1}+q_{m-2}\longrightarrow\infty,
\]
one has $Q_m>S_m>0$ for every sufficiently large $m$.  The Euclidean
algorithm on the coprime pair $(Q_m,S_m)$ constructs its positive
continued-fraction denominator word.  The two finite presentations obtained
by splitting the last partial quotient choose the required determinant sign.
After that choice, any two numerator rows with the same lower row and the
same determinant differ by an integral multiple of $(Q_m,S_m)$; changing
the first partial quotient absorbs precisely that multiple.  Thus the
composite matrix is a finite continued-fraction prefix, and the two numbers
have a common tail.  This proves the needed Serret direction without the
printed circular step.  The complete locator audit is in
\path{qa/STEUDING-2005-F011-exercise-522-audit.md}.
\end{sourceissue}

\begin{theorem}[edition solution of Steuding's Exercise 5.22]
\label{thm:Steuding-522-equivalence}
Let $b_{3r}=4$, $b_{3r+1}=2$, and $b_{3r+2}=1$.  The following statements
are equivalent:
\begin{align*}
 &\text{for every $m\in\Npos$, the $U$-orbit of $m$ eventually enters the
 cycle $4,2,1$};\\
 &\text{for every $m\in\Npos$, }\Theta(m)\sim[\overline{4,2,1}].
\end{align*}
More precisely, for each fixed $m$,
\begin{equation}
\label{eq:Steuding-pointwise-equivalence}
 \Theta(m)\sim[\overline{4,2,1}]
 \quad\Longleftrightarrow\quad
 \exists r,s\in\Nzero\ \forall j\in\Nzero:
 U^{r+j}(m)=b_{s+j}.
\end{equation}
\end{theorem}

\begin{proof}
If the orbit eventually enters the cycle, its digit sequence in
\eqref{eq:Steuding-Theta} is eventually a cyclic shift of
$(4,2,1,4,2,1,\ldots)$.  Serret's theorem therefore gives the forward
implication of \eqref{eq:Steuding-pointwise-equivalence}.  Conversely,
Serret's theorem turns the equivalence of the two irrational continued
fractions into the displayed equality of tails for some $r,s$.  Hence
$U^r(m)\in\{4,2,1\}$.  The literal identities
\[
 U(4)=2,
 \qquad U(2)=1,
 \qquad U(1)=4
\]
then force every subsequent value.  Universal quantification over $m$ gives
the two global statements.  The phase variable $s$ is retained; no preferred
entry point into the cycle is assumed.
\end{proof}

\begin{calculation}[the requested Markoff constant]
\label{calc:Steuding-522-Markoff}
Let $d=\sqrt{229}$.  The exact period matrix is
\[
 \begin{pmatrix}4&1\\1&0\end{pmatrix}
 \begin{pmatrix}2&1\\1&0\end{pmatrix}
 \begin{pmatrix}1&1\\1&0\end{pmatrix}
 =\begin{pmatrix}13&9\\3&2\end{pmatrix}.
\]
Therefore
\begin{equation}
\label{eq:Steuding-beta-quadratic}
 3\beta^2-11\beta-9=0,
 \qquad
 \beta=\frac{11+d}{6}.
\end{equation}
Index its partial quotients by $a_{3r}=4$, $a_{3r+1}=2$,
$a_{3r+2}=1$.  Formula~(5.11) of the source is
\begin{equation}
\label{eq:Steuding-Markoff-Perron}
 \lambda(\beta)=\liminf_{n\to\infty}
 \left(
  [a_{n+1},a_{n+2},\ldots]
  +[0,a_n,a_{n-1},\ldots,a_1]
 \right)^{-1}.
\end{equation}
The three phases, with no cyclic normalization, are
\[
\begin{array}{c|c|c|c|c}
n\bmod3 & [a_{n+1},a_{n+2},\ldots]
 & \displaystyle\lim[0,a_n,\ldots,a_1]
 & \text{sum} & \text{reciprocal}\\ \hline
0 & (13+d)/10 & (d-13)/10 & d/5 & 5/d\\
1 & (7+d)/18  & (d-7)/18  & d/9 & 9/d\\
2 & (11+d)/6  & (d-11)/6  & d/3 & 3/d
\end{array}
\]
The reversed finite fractions converge on each residue subsequence to the
displayed reversed periodic tail; equivalently, Steuding's Galois formula
identifies each with the negative conjugate of the corresponding cyclic
complete quotient \cite[Thm.~5.3 and Eq.~(5.11)]{Steuding2005}.  Taking the
liminf in \eqref{eq:Steuding-Markoff-Perron} gives
\begin{equation}
\label{eq:Steuding-Markoff-value}
 \boxed{\lambda([\overline{4,2,1}])
 =\min\left\{\frac5{\sqrt{229}},\frac9{\sqrt{229}},
                    \frac3{\sqrt{229}}\right\}
 =\frac3{\sqrt{229}}.}
\end{equation}
The exact quadratic-field identities are reproduced by
\path{certificates/steuding_522_checks.py}.
\end{calculation}

\begin{proposition}[exact Crandall clock crosswalk]
\label{prop:Steuding-Crandall-crosswalk}
For an odd start $y_0$, put $y_i=C_{\mathrm{Cr}}^i(y_0)$,
$e_i=\vtwo(3y_i+1)$, and $A_j=\sum_{i<j}e_i$.  Then the Steuding digit
sequence and the accelerated orbit satisfy
\begin{equation}
\label{eq:Steuding-Crandall-clock}
 y_j=U^{A_j+j}(y_0),
\end{equation}
and, for $1\leq r\leq e_j+1$,
\begin{equation}
\label{eq:Steuding-inserted-even-states}
 U^{A_j+j+r}(y_0)=\frac{3y_j+1}{2^{r-1}}.
\end{equation}
For an arbitrary $m=2^s y_0$ with $y_0$ odd, the first $s$ Steuding digits
are the exact halving prefix before this odd-return expansion begins.
\end{proposition}

\begin{proof}
Equation \eqref{eq:Steuding-Crandall-clock} is
\eqref{eq:Crandall-three-clocks}.  Starting from the odd state $y_j$, the
first $U$ step forms $3y_j+1$ and the next $e_j$ steps divide successively by
$2$, proving \eqref{eq:Steuding-inserted-even-states}; at $r=e_j+1$ the
value is $y_{j+1}$.  The even-start statement is the literal iteration
$U^r(2^s y_0)=2^{s-r}y_0$ for $0\leq r\leq s$.
\end{proof}

\begin{nonclaim}
\label{nonclaim:Steuding-Markoff-boundary}
Theorem~\ref{thm:Steuding-522-equivalence} is a lossless endpoint recoding,
not a proof of either side: the full orbit is stored as the full digit
sequence and the initial integer is recovered by the floor map.  Steuding's
Theorem~5.8 gives only
\[
 \Theta(m)\sim\beta
 \ \Longrightarrow\ 
 \lambda(\Theta(m))=\frac3{\sqrt{229}}.
\]
Equality of Markoff constants is not a converse and cannot certify Collatz
termination.  The orbit-digit continued fraction $\Theta(m)$ is also not
Crandall's Section~7 continued fraction for the fixed number $\log_2 3$;
no map between those two constructions is evidenced here.  Finally,
\eqref{eq:Steuding-Crandall-clock} is a variable-time expansion, not a
one-step identification of $U$ and $C_{\mathrm{Cr}}$.
\end{nonclaim}
